\scanpage{193}
\mspage{II.41}
% archive sheet number: 192
% block: top, 6.11 | display
\[
6.11.\qquad \operatorname{Mod}_{\mathbb{Q}_\ell}(\Pi^{\mathrm{an}})\simeq
\operatorname{Mod}(\Pi^{\mathrm{alg}})\simeq
\varinjlim_\lambda\operatorname{Mod}_{\mathbb{Q}_\ell}(\pi_\lambda)
\]
% block: upper third, 6.12 | text
6.12. Un cas \unsure{important} est celui où dans le \deleted{système}
pro-objet $\Pi=(\pi_\lambda)$, les morphismes de transition sont
\unsure{quasi-surjectifs} (i.e. image d'indice fini). Alors il en est de même
dans $\Pi^{\mathrm{an}}$ et $\Pi^{\mathrm{alg}}$ ; (en particulier
$(\Pi^{\mathrm{alg}})^0$ est un pro-objet \uline{strict}). De plus, les syst.
projectifs $\mathfrak{g}=\operatorname{Lie}(\Pi^{\mathrm{an}})$ et
$\widetilde{\mathfrak{g}}=\operatorname{Lie}(\Pi^{\mathrm{alg}})$ sont aussi
stricts \unsure{....|...}

% block: middle third, 6.13 | text
6.13. Soit $X=\varprojlim_\lambda X_\lambda$ un schéma, limite projective
d'un système projectif filtrant \inserted{$\unsure{\mathscr{X}|\mathcal{X}}$}
de schémas, à morphismes de transition affines. Soit $\unsure{\xi|\zeta}$ un pt géom.
de $X$, $\xi_\lambda$ son image dans $X_\lambda$. Considérons
% block: middle, fundamental pro-group | display
\[
\Pi=\pi_1(\mathscr{X},\xi)=(\pi_1(X_\lambda,\xi))_\lambda,
\]
% block: lower third, representations | text
système projectif de groupes profinis. Les représentations
$\ell$-adiques de ce \deleted{système} pro-groupe top. forment une
catégorie équivalente à la catégorie
$\varinjlim_\lambda(\inserted{\mathbb{Q}_\ell\text{-}}\text{Faisceaux
\deleted{constants} tordus sur }X_\lambda)$, d'où leur intérêt. Cette
catégorie s'explicite en termes de $\Pi^{\mathrm{an}}$ ou $\Pi^{\mathrm{alg}}$
comme on a vu plus haut, et peut aussi (partiellement) s'étudier à l'aide
de $\operatorname{Lie}(\Pi^{\mathrm{an}})$ et
$\operatorname{Lie}(\Pi^{\mathrm{alg}})$
